\section*{Problem 11}

Use the generalized Chinese Remainder Theorem.

\subsection*{(a)}
Solve
\[
x \equiv 5 \pmod{6},\qquad x \equiv 3 \pmod{10},\qquad x \equiv 8 \pmod{15}.
\]
Consistency checks:
\[
\begin{aligned}
&\gcd(6,10)=2: &&5\equiv 3 \equiv 1 \pmod{2}\quad \checkmark\\
&\gcd(6,15)=3: &&5\equiv 8 \equiv 2 \pmod{3}\quad \checkmark\\
&\gcd(10,15)=5:&&3\equiv 8 \equiv 3 \pmod{5}\quad \checkmark
\end{aligned}
\]
So a solution exists and is unique modulo $\mathrm{lcm}(6,10,15)=30$.

From $x\equiv 5 \pmod{6}$ the candidates modulo $30$ are $5,11,17,23,29$.  
Among these, those $\equiv 3 \pmod{10}$ are $23$. Check $23\equiv 8 \pmod{15}$, true.

\[
\boxed{x \equiv 23 \pmod{30}}.
\]

\subsection*{(b)}
Solve
\[
x \equiv 7 \pmod{9},\qquad x \equiv 4 \pmod{12},\qquad x \equiv 16 \pmod{21}.
\]
Consistency checks (all use $\gcd=3$): $7\equiv4\equiv16\equiv 1 \pmod{3}$, so consistent.
Uniqueness is modulo $\mathrm{lcm}(9,12,21)=252$.

First combine mod $9$ and $12$.
Write $x=7+9a$. Then $7+9a\equiv 4 \pmod{12}\iff 9a\equiv 9 \pmod{12}$,
which reduces to $3a\equiv 3 \pmod{4}\iff a\equiv 1 \pmod{4}$.
Thus $x\equiv 16 \pmod{36}$.

Now impose $x\equiv 16 \pmod{21}$ as well:
$x=16+36k \equiv 16+15k \pmod{21}$, so $15k\equiv 0 \pmod{21}$,
which gives $k\equiv 0 \pmod{7}$.
Hence $x=16+252t$.

\[
\boxed{x \equiv 16 \pmod{252}}.
\]
